\documentclass[10pt,a4paper]{amsart}
\usepackage{amsmath,amssymb,amsthm,graphicx}
% Retained minimal preamble for the Stacks Project chapter flat.tex (More on Flatness).
% Verbatim excerpt of the chapter preamble: Xy-pic setup (native preamble lines 14-18) and the native
% theorem declarations with the mathematical macros (native preamble lines 154-217). Copied mechanically;
% nothing added, removed or redefined. The wrapper class and the amsmath/amssymb/amsthm/graphicx base
% packages are supplied by the assembler, so the class/externaldocument part of the native preamble is omitted.
% For commutative diagrams we use Xy-pic
\usepackage[all]{xy}

% We use 2cell for 2-commutative diagrams.
\xyoption{2cell}
% Theorem environments.
%
\theoremstyle{plain}
\newtheorem{theorem}[subsection]{Theorem}
\newtheorem{proposition}[subsection]{Proposition}
\newtheorem{lemma}[subsection]{Lemma}

\theoremstyle{definition}
\newtheorem{definition}[subsection]{Definition}
\newtheorem{example}[subsection]{Example}
\newtheorem{exercise}[subsection]{Exercise}
\newtheorem{situation}[subsection]{Situation}

\theoremstyle{remark}
\newtheorem{remark}[subsection]{Remark}
\newtheorem{remarks}[subsection]{Remarks}

\numberwithin{equation}{subsection}

% Macros
%
\def\lim{\mathop{\mathrm{lim}}\nolimits}
\def\colim{\mathop{\mathrm{colim}}\nolimits}
\def\Spec{\mathop{\mathrm{Spec}}}
\def\Hom{\mathop{\mathrm{Hom}}\nolimits}
\def\Ext{\mathop{\mathrm{Ext}}\nolimits}
\def\SheafHom{\mathop{\mathcal{H}\!\mathit{om}}\nolimits}
\def\SheafExt{\mathop{\mathcal{E}\!\mathit{xt}}\nolimits}
\def\Sch{\mathit{Sch}}
\def\Mor{\mathop{\mathrm{Mor}}\nolimits}
\def\Ob{\mathop{\mathrm{Ob}}\nolimits}
\def\Sh{\mathop{\mathit{Sh}}\nolimits}
\def\NL{\mathop{N\!L}\nolimits}
\def\CH{\mathop{\mathrm{CH}}\nolimits}
\def\proetale{{pro\text{-}\acute{e}tale}}
\def\etale{{\acute{e}tale}}
\def\QCoh{\mathit{QCoh}}
\def\Ker{\mathop{\mathrm{Ker}}}
\def\Im{\mathop{\mathrm{Im}}}
\def\Coker{\mathop{\mathrm{Coker}}}
\def\Coim{\mathop{\mathrm{Coim}}}

% Boxtimes
%
\DeclareMathSymbol{\boxtimes}{\mathbin}{AMSa}{"02}

%
% Macros for moduli stacks/spaces
%
\def\QCohstack{\mathcal{QC}\!\mathit{oh}}
\def\Cohstack{\mathcal{C}\!\mathit{oh}}
\def\Spacesstack{\mathcal{S}\!\mathit{paces}}
\def\Quotfunctor{\mathrm{Quot}}
\def\Hilbfunctor{\mathrm{Hilb}}
\def\Curvesstack{\mathcal{C}\!\mathit{urves}}
\def\Polarizedstack{\mathcal{P}\!\mathit{olarized}}
\def\Complexesstack{\mathcal{C}\!\mathit{omplexes}}
% \Pic is the operator that assigns to X its picard group, usage \Pic(X)
% \Picardstack_{X/B} denotes the Picard stack of X over B
% \Picardfunctor_{X/B} denotes the Picard functor of X over B
\def\Pic{\mathop{\mathrm{Pic}}\nolimits}
\def\Picardstack{\mathcal{P}\!\mathit{ic}}
\def\Picardfunctor{\mathrm{Pic}}
\def\Deformationcategory{\mathcal{D}\!\mathit{ef}}

% --- appended 2026-10-12 (Chief): the project own definitions of the reference/slogan/history
% display environments, copied mechanically from the shared Stacks preamble lines 8-12
% (\usepackage{verbatim} plus the three \newenvironment lines, which the chapter preamble keeps
% verbatim from the project preamble). Nothing added, removed or redefined: these environments
% carry no mathematical content and the project itself typesets them as comments. Without these
% three lines a counted statement whose own text uses them fails with Environment slogan undefined.
\usepackage{verbatim}
\newenvironment{reference}{\comment}{\endcomment}
\newenvironment{slogan}{\comment}{\endcomment}
\newenvironment{history}{\comment}{\endcomment}
\usepackage{cleveref}
\crefname{originalequation}{original equation}{original equations}
\Crefname{originalequation}{Original equation}{Original equations}
\makeatletter
\newlabel{equation-star}{{0AT6}{source}{}{sourcepointer-1}{}}
\newlabel{equation-star@cref}{{[originalequation][0][]0AT6}{[1][0][]source}{}{}{}}
\makeatother
\crefname{originallemma}{original lemma}{original lemmas}
\Crefname{originallemma}{Original lemma}{Original lemmas}
\makeatletter
\newlabel{lemma-invert-universally-injective}{{05FR}{source}{}{sourcepointer-2}{}}
\newlabel{lemma-invert-universally-injective@cref}{{[originallemma][0][]05FR}{[1][0][]source}{}{}{}}
\makeatother
\title{Existence of a finitely generated ideal over which a finitely presented module has a pure spreadout}
\author{The Stacks Project Authors}
\date{stacks-project:flat.tex:local-snapshot-20261011; GNU FDL 1.2 or later}
\begin{document}
\maketitle
\noindent\textit{Editorial note.} All delivered mathematical text is the author's own, copied from the pinned Stacks Project chapter file flat.tex; no statement or proof step was rewritten, reordered or removed.\par\medskip
\noindent\textit{Editorial note.} The preamble is the chapter's own retained minimal preamble (sources/stacks-flat/retained-preamble.tex), a verbatim excerpt of the native Xy-pic setup and of the native theorem declarations and mathematical macros, followed by the project own definitions of the reference, slogan and history display environments (verbatim package plus three newenvironment lines) that the retained file carries at its end and that the project itself typesets as comments. No macro was invented or redefined, and the wrapper class and the amsmath/amssymb/amsthm/graphicx base packages are supplied by the assembler.\par\medskip
\noindent\textit{Editorial note.} The counted statement and proof contain no cross reference, no citation and no included file, so the delivered body is self contained.\par\medskip
\noindent\textit{Editorial note.} Canonical location of the counted proposition: Stacks tag 0AT8 (sources/stacks-flat/pinned-original-source.tex line 7270).\par\medskip
\noindent\textit{Allowed standard prerequisites (editorial).}\begin{itemize}
\item The situation of the reference equation and the notion of a pure spreadout. Same chapter, flat.tex, Equation labelled equation-star (source line 7208), invoked in the statement. Fixes the rings A, B, the localisation S and the module N in which the spreadout is sought.
\item Universal injectivity of a map of modules with a free target after dividing out. Same chapter, Lemma labelled lemma-invert-universally-injective (source line 1976), invoked at proof lines 7298-7299. Supplies that the injectivity of multiplication by the elements of S on N'/mN' implies universal injectivity of N'/IN' -> N/IN; the construction of I and of the finite model is inside the counted span.
\end{itemize}\medskip
\ifdefined\hypertarget\hypertarget{sourcepointer-1}{}\fi
\noindent\textit{Source reference.} Original equation 0AT6: sources/stacks-flat/pinned-original-source.tex, source lines 7208--7208. Same-chapter pointer into the author environment equation at source line 7208; invoked at proof line None of the counted proof. It is retained only as a navigation pointer; every load-bearing step is inside the delivered span.\par\medskip
\ifdefined\hypertarget\hypertarget{sourcepointer-2}{}\fi
\noindent\textit{Source reference.} Original lemma 05FR: sources/stacks-flat/pinned-original-source.tex, source lines 1976--1976. Same-chapter pointer into the author environment lemma at source line 1976; invoked at proof line 7299 of the counted proof. It is retained only as a navigation pointer; every load-bearing step is inside the delivered span.\par\medskip
\begin{lemma}
\label{lemma-find-pure-spreadout}
In (\ref{equation-star}) for every $\mathfrak p \in \Spec(A)$
there is a finitely generated ideal $I \subset \mathfrak pA_\mathfrak p$
such that over $A_\mathfrak p/I$ we have a pure spreadout.
\end{lemma}

\begin{proof}
We may replace $A$ by $A_\mathfrak p$. Thus we may assume $A$ is
local and $\mathfrak p$ is the maximal ideal $\mathfrak m$ of $A$.
We may write $N = S^{-1}N'$ for some finitely presented $B$-module $N'$
by clearing denominators in a presentation of $N$ over $S^{-1}B$.
Since $B/\mathfrak m B$ is Noetherian, the kernel $K$ of
$N'/\mathfrak m N' \to N/\mathfrak m N$ is finitely generated.
Thus we can pick $s \in S$ such that $K$ is annihilated by $s$.
After replacing $B$ by $B_s$ which is allowed as it just means passing
to an affine open subscheme of $\Spec(B)$, we find that the elements of $S$
are injective on $N'/\mathfrak m N'$. At this point we choose
a local subring $A_0 \subset A$ essentially of finite type over $\mathbf{Z}$,
a finite type ring map $A_0 \to B_0$ such that $B = A \otimes_{A_0} B_0$,
and a finite $B_0$-module $N'_0$ such that
$N' = B \otimes_{B_0} N'_0 = A \otimes_{A_0} N'_0$.
We claim that $I = \mathfrak m_{A_0} A$ works.
Namely, we have
$$
N'/IN' = N'_0/\mathfrak m_{A_0} N'_0 \otimes_{\kappa_{A_0}} A/I
$$
which is free over $A/I$. Multiplication by the elements of $S$
is injective after dividing out by the maximal ideal, hence
$N'/IN' \to N/IN$ is universally injective for example by
Lemma \ref{lemma-invert-universally-injective}.
\end{proof}

\end{document}
